\chapter{Style smoke test}\label{ch:smoke}
\index{style}
This chapter exists only to prove that every environment, figure type and code block compiles. Delete it when the first real chapter exists. The model for structure is \cite{sannemo2025}.

\section{Environments}
\begin{definition}[Function]\index{function}
A \emph{function} $f:A\to B$ assigns to each $a\in A$ exactly one $f(a)\in B$.
\end{definition}
\begin{theorem}[Sum of the first $n$ integers]
For every $n\ge 1$, $\sum_{k=1}^{n}k=\frac{n(n+1)}{2}$.
\end{theorem}
\begin{proof}
By induction on $n$. For $n=1$ both sides are 1. If it holds for $n$, then adding $n+1$ gives $\frac{n(n+1)}2+n+1=\frac{(n+1)(n+2)}2$.
\end{proof}
\begin{lemma}A small lemma.\end{lemma}
\begin{proposition}A small proposition.\end{proposition}
\begin{example}\label{ex:one}
Worked example: strategy first, then solution.\par
\begin{solution}The answer is $42$.\end{solution}
\end{example}
\begin{remark}A remark.\end{remark}
\begin{exercise}[easy] Show that $\sum_{k=1}^{n}(2k-1)=n^2$.\end{exercise}
\begin{exercise}[hard] Prove the theorem above without induction.\end{exercise}
\begin{principle}[Principle]A principle box.\end{principle}
\begin{rulebox}[Decision rule]A rule box.\end{rulebox}
\begin{tripwire}[Tripwire]A warning box.\end{tripwire}
\begin{wins}[Definition of done]A done box.\end{wins}
\begin{checkpoint}[Try it] Run the code in Listing~\ref{lst:hello} and predict the output first.\end{checkpoint}

\section{A graph}
\begin{figure}[ht]\centering
\begin{tikzpicture}
\begin{axis}[width=0.8\linewidth,height=5.5cm,xlabel={$x$},ylabel={$f(x)$},legend pos=north west,grid=major,
  grid style={gray!25}]
\addplot[domain=0:4,samples=40,thick,black]{x^2};
\addplot[domain=0:4,samples=40,thick,dashed,accent]{2*x};
\legend{$x^2$,$2x$}
\end{axis}
\end{tikzpicture}
\caption{A pgfplots graph. Notice that $x^2$ overtakes $2x$ after $x=2$. Line styles, not only colour, tell the curves apart.}\label{fig:plot}
\end{figure}

\section{A diagram, step by step}
\begin{figure}[ht]\centering
\foreach \s in {1,2,3}{%
\begin{tikzpicture}[scale=0.8,every node/.style={circle,draw,minimum size=6mm,inner sep=0pt,font=\small}]
\node (a) at (0,0) {A};\node (b) at (1.2,0.8) {B};\node (c) at (1.2,-0.8) {C};\node[fill=accentsoft] (d) at (2.4,0) {D};
\draw (a)--(b) (a)--(c) (b)--(d) (c)--(d);
\node[draw=none,rectangle,font=\footnotesize] at (1.2,-1.6) {step \s};
\end{tikzpicture}\quad}
\caption{Small multiples show the state of an algorithm at three successive steps.}\label{fig:steps}
\end{figure}

\section{A circuit}
\begin{figure}[ht]\centering
\begin{circuitikz}[scale=0.9,transform shape]
\draw (0,0) to[battery1,l=$V$] (0,2) to[R,l=$R$] (3,2) to[C,l=$C$] (3,0) -- (0,0);
\end{circuitikz}
\caption{An RC circuit drawn with circuitikz. Every circuit in the book is drawn this way.}\label{fig:rc}
\end{figure}

\section{Code with output}
\begin{lstlisting}[language=Python,caption={A short listing with line numbers.},label={lst:hello}]
# why: show the loop invariant, not just the loop
total = 0
for k in range(1, 6):
    total += k
print(total)
\end{lstlisting}
Output of a real run:
\begin{lstlisting}[numbers=none]
15
\end{lstlisting}

\begin{table}[ht]\centering
\begin{tabularx}{\linewidth}{@{}lYr@{}}
\toprule
\tablehead\thead{Name} & \thead{Meaning} & \thead{Cost}\\
\midrule
Sort & Order the items & $O(n\log n)$\\
Scan & Visit each item once & $O(n)$\\
\bottomrule
\end{tabularx}
\caption{Cheat sheet table: each chapter ends its main text with one of these.}
\end{table}

\section*{Notes}
\addcontentsline{toc}{section}{Notes}
Sources are cited by number~\cite{sannemo2025}.
